\subsection{Agent-oriented software engineering}
\label{sec:aose}
\slidesrc{AOSE: 1--8}
\movedfrom{the AOSE deck, shown right after the EMAS keynote}

The short AOSE deck (based on Zambonelli and Omicini, \emph{Challenges and Research
Directions in Agent-Oriented Software Engineering}, 2004, and Giorgini and
Henderson-Sellers, \emph{Agent-Oriented Methodologies: An Introduction}, 2005) explains
\emph{why} agents are interesting for software engineering at all.

\begin{nbfact}{Why agents suit modern software (slides 4--5)}
The features of agent-based systems are well suited to tackle the complexity of developing
software in modern scenarios:
\begin{enumerate}
  \item the \textbf{autonomy} of application components reflects the intrinsically
        \emph{decentralised} nature of modern distributed systems, and can be considered the
        natural extension of the notions of \emph{modularity} and \emph{encapsulation} to
        systems that are owned by different stakeholders;
  \item the \textbf{flexible} way in which agents operate and interact (with each other and
        with the environment) suits the \emph{dynamic and unpredictable} scenarios where
        software is expected to operate;
  \item the concept of \textbf{agency} provides a unified view of AI results and
        achievements, making agents and MAS act as sound and manageable
        \emph{repositories of intelligent behaviours}.
\end{enumerate}
\end{nbfact}

\begin{nbintuition}[Autonomy as ``encapsulation of behaviour'']
Object orientation encapsulates \emph{state}: nobody outside can change an object's private
fields. Agent orientation goes one step further and encapsulates \emph{behaviour}: nobody
outside can force an agent to act; you can only ask (Section~\ref{sec:agents-history}). When
the components belong to different companies (your bank's agent, a shop's agent, a
courier's agent), this is not a design choice but a fact of life.
\end{nbintuition}

\begin{nbdef}{Agent-Oriented Software Engineering (slide 6)}
With the increasing acceptance of agent-based computing as a novel software engineering
paradigm, there has been a great deal of research on suitable models and techniques to
support the development of complex software systems in terms of MAS. This research is
roughly grouped under the term \textbf{Agent-Oriented Software Engineering} (AOSE).
\end{nbdef}

\begin{nbfact}{Mainstream research directions in AOSE (AOSE slide 8; EMAS slide 11)}
\begin{center}
\small
\begin{tabularx}{\linewidth}{@{}c>{\raggedright\arraybackslash}X|c>{\raggedright\arraybackslash}X@{}}
\toprule
\multicolumn{2}{c|}{\textbf{AOSE deck (2004 view)}} & \multicolumn{2}{c}{\textbf{EMAS keynote (2024 view)}} \\
\midrule
1 & agent-oriented architectures & 1 & agent-oriented architectures \\
2 & AOSE methodologies & 2 & methodologies and notation techniques \\
3 & notation techniques & 3 & agent-oriented programming languages \\
4 & agent-oriented programming languages & 4 & MAS infrastructures and toolkits \\
5 & agent-based modelling and simulation & 5 & agent-based modelling and simulation \\
6 & MAS infrastructures & 6 & symbolic and sub-symbolic integration \\
  & & 7 & engineering safe autonomous systems \\
\bottomrule
\end{tabularx}
\end{center}
Twenty years later the list is almost the same, plus two new directions: neurosymbolic
integration and safety.
\end{nbfact}
(A section of the original AOSE deck, ``The promise of AOSE for distributed systems
engineering'', was removed in the shortened version used in class.)
